% ============================================================================
% 09_modelo_datos.tex — Cuerpo del entregable "Modelo de Base de Datos"
% Sin preámbulo: se incluye tanto desde 09_Modelo_Datos.tex (standalone)
% como, eventualmente, desde un documento maestro.
% ============================================================================
\section{Modelo de Base de Datos}

\subsection{Justificación de la persistencia basada en archivos}

LSA Salud no utiliza un motor de base de datos relacional. Toda la información
que el sistema persiste se guarda como archivos en el sistema de archivos
local del backend (\texttt{dataset/}, \texttt{models/}, \texttt{voices/}).
Esta decisión de diseño responde a las siguientes razones, verificables
contra el código y contra la ERS de Fase 1:

\begin{itemize}
    \item \textbf{Escala de despliegue.} RNF-T-02 establece que cada
    estación fija opera de forma autónoma en su propia PC local, con un
    par paciente-personal por instancia de backend y sin datos
    compartidos con las demás estaciones del sistema (admisión, recepción,
    triaje, guardia). No existe un escenario de múltiples instancias
    concurrentes escribiendo sobre el mismo dataset, por lo que los
    problemas que una base de datos relacional resuelve (control de
    concurrencia entre muchos escritores, replicación, transacciones
    distribuidas) no se presentan en este contexto.

    \item \textbf{Formato nativo del modelo.} Las muestras de entrenamiento
    se capturan y se consumen como arreglos \texttt{numpy}
    (\texttt{float32}) y los pesos del modelo como \texttt{state\_dict} de
    PyTorch. Guardarlos directamente como \texttt{.npy} y \texttt{.pth} evita
    una capa de serialización/deserialización adicional (por ejemplo, un
    ORM o un mapeo objeto-relacional) entre el formato de disco y el
    formato que \texttt{train\_lstm.py} y \texttt{app.py} usan en memoria.

    \item \textbf{Sin necesidad de consultas relacionales.} El acceso a los
    datos persistidos es siempre por lectura completa de un directorio
    (todas las muestras de una etiqueta, para \texttt{GET /dataset\_stats}
    o para \texttt{train\_lstm.py}) o por nombre de archivo exacto (una
    muestra puntual). No hay filtros, \emph{joins} ni agregaciones más
    complejas que un conteo de archivos por carpeta, por lo que el costo de
    modelar y mantener un esquema relacional no se justifica.

    \item \textbf{Privacidad y localidad.} RNF-03.1 exige que todo el
    procesamiento y almacenamiento ocurra en la máquina local, sin que el
    dato de salud del paciente salga del equipo. Un archivo en el sistema de
    archivos local ya cumple ese objetivo sin necesidad de un servidor de
    base de datos adicional (con su propia superficie de configuración,
    puertos y credenciales) que administrar.

    \item \textbf{La conversación no se persiste, por decisión de
    producto.} El historial de mensajes entre paciente y personal de
    recepción o enfermería
    vive únicamente como estado de React sincronizado entre pestañas por
    \texttt{BroadcastChannel} (ver \texttt{web/hooks/useChatChannel.ts}) y
    se pierde al cerrar o recargar la pestaña. Esto no es una limitación
    técnica: es consistente con RNF-03.1 (no persistencia de conversaciones)
    y evita que LSA Salud acumule un historial de conversaciones médicas del
    paciente sin necesidad funcional de conservarlo.
\end{itemize}

En consecuencia, este entregable no presenta un esquema de base de datos
relacional, sino un modelo de los datos que LSA Salud persiste, expresado con
las mismas técnicas de modelado (entidad-relación) que se usarían para una
base de datos, pero mapeado explícitamente a archivos y directorios en la
Sección~\ref{sec:dat-mapeo-fisico}.

\subsection{Modelo conceptual (Entidad-Relación)}

La Figura~\ref{fig:dat-er-conceptual} presenta el modelo conceptual de la
información persistida por LSA Salud. Se reutilizan, con el mismo nombre y
significado, los conceptos \textit{ClaseDeSeña}, \textit{MuestraDeEntrenamiento}
y \textit{ModeloDeReconocimiento} ya definidos en el Modelo de Dominio
(Entregable 01); a diferencia de aquel diagrama —que modela el problema sin
elementos de diseño—, aquí estos mismos conceptos se presentan enfocados en
lo que efectivamente queda grabado en disco.

\diagrama{09_er_conceptual}{\textwidth}{Diagrama Entidad-Relación conceptual de la persistencia de LSA Salud}{fig:dat-er-conceptual}

Se agregan dos conceptos sin equivalente directo en el Modelo de Dominio,
porque pertenecen al diseño de persistencia y no al problema de negocio:
\textit{VocabularioDelModelo} (la lista ordenada de clases que un
\textit{ModeloDeReconocimiento} concreto es capaz de reconocer) y
\textit{ArchivoDeVoz} (el modelo de síntesis de voz Piper, un activo de
terceros pre-entrenado, sin relación con el resto de las entidades). Cada
\textit{MuestraDeEntrenamiento} persiste, en un mismo registro, tanto los
landmarks de manos y pose como los blendshapes faciales del gesto: MediaPipe
extrae ambas señales en una única pasada por frame, por lo que no existe una
entidad de blendshapes separada. Quedan fuera de este modelo, por no
persistirse, los conceptos \textit{Conversación}, \textit{Mensaje} y los
cachés de traducción y de audio; se documentan igualmente en la
Sección~\ref{sec:dat-mapeo-fisico} para dejar constancia explícita de que
fueron considerados y descartados.

Se agregan además dos entidades que cubren los requisitos RF-03,
RF-04 y RF-06: \textit{EventoDeAuditoría} (RF-04.1, cuyo campo \texttt{tipo}
también cubre el consentimiento de RF-03.1 con el valor \texttt{consent} ---
no existe una entidad \textit{EventoDeConsentimiento} separada ni un
archivo \texttt{logs/consent.jsonl}: el consentimiento es un evento de
auditoría más) y \textit{ReporteDeTraducción} (RF-06.1). Ambas son, por
diseño, anónimas --- sin identidad del paciente ni contenido de imagen o
video --- consistentes con RNF-03.1, y se persisten con el mismo patrón
basado en archivos que el resto del modelo (ver
Sección~\ref{sec:dat-mapeo-fisico}).

\subsection{Modelo lógico (notación pata de gallo)}

La Figura~\ref{fig:dat-modelo-logico} traduce el modelo conceptual a un
esquema lógico con claves primarias, claves foráneas y cardinalidad en
notación pata de gallo (\emph{crow's foot}), tal como se usaría para diseñar
las tablas de una base de datos relacional.

\diagrama{09_modelo_logico}{\textwidth}{Modelo lógico de la persistencia de LSA Salud (notación pata de gallo)}{fig:dat-modelo-logico}

\textbf{Aclaración importante:} estas ``entidades'' no son tablas de una
base de datos. Cada una es, en realidad, un patrón de archivo o de
directorio del sistema de archivos local:

\begin{itemize}
    \item La clave primaria compuesta de \textit{MuestraDeEntrenamiento}
    (\texttt{etiqueta} + \texttt{timestamp\_ms}) se materializa como la ruta
    física \texttt{dataset/\allowbreak{}custom/\allowbreak{}<etiqueta>/\allowbreak{}<timestamp\_ms>.npy}
    para los landmarks, con los blendshapes del mismo gesto en el archivo
    hermano
    \texttt{dataset/\allowbreak{}custom/\allowbreak{}<etiqueta>/\allowbreak{}<timestamp\_ms>.blendshapes.npy}:
    \texttt{etiqueta} es el nombre del subdirectorio y \texttt{timestamp\_ms}
    es el nombre base compartido por ambos archivos, emparejados por
    convención de nombre y no por una clave foránea almacenada en ningún
    lado.
    \item \textit{ClaseDeSeña} no tiene un archivo propio: es un valor
    derivado (el nombre de cada subdirectorio bajo \texttt{dataset/custom/}).
    \item \textit{VocabularioDelModelo} es el arreglo JSON completo de
    \texttt{models/custom\_classes.json}; el ``índice'' de cada fila es la
    posición dentro del arreglo, que a su vez es el índice de salida de la
    capa final del modelo.
    \item \textit{ModeloDeReconocimiento} es un único archivo binario
    (\texttt{state\_dict} de PyTorch), no una fila de una tabla.
\end{itemize}

Ningún motor de base de datos participa en la lectura o escritura de estas
``entidades''; toda la integridad referencial (por ejemplo, que
\texttt{VocabularioDelModelo} esté alineado con la capa de salida del
\texttt{ModeloDeReconocimiento} cargado) se garantiza por convención en el
código (\texttt{app.py}, \texttt{train\_lstm.py}) y no por restricciones de
esquema.

\subsection{Diccionario de datos}

Las Tablas~\ref{tab:dat-dic-muestra} a~\ref{tab:dat-dic-conversacion}
detallan, para cada entidad lógica, sus atributos, tipos, formato/forma y
restricciones. Las restricciones de validación provienen de
\texttt{capture\_validation.py} y de las verificaciones adicionales en el
endpoint \texttt{POST /capture} de \texttt{app.py}.

\begin{longtable}{p{2.8cm} p{2.2cm} p{3.2cm} p{5.6cm}}
    \toprule
    \textbf{Atributo} & \textbf{Tipo} & \textbf{Formato / forma} & \textbf{Restricciones} \\
    \midrule
    \endhead
    etiqueta & \texttt{str} & directorio & Debe cumplir la expresión regular
    \texttt{\textasciicircum{}[a-z0-9\_]\{1,32\}\$} (validado por
    \texttt{validate\_capture\_label}). \\

    timestamp\_ms & \texttt{int64} & epoch ms (nombre de archivo) &
    Generado por el backend como \texttt{int(time.time() * 1000)}; único
    dentro de la carpeta de la etiqueta (dos capturas en el mismo
    milisegundo sobrescribirían el archivo previo); es el nombre base
    compartido por el archivo de landmarks y el de blendshapes de la misma
    muestra. \\

    frames & \texttt{float32} & \texttt{(45, 75, 4)} & El backend exige un
    mínimo de 5 frames por muestra (\texttt{len(frames) < 5} rechaza la
    petición); el flujo de captura normal siempre envía 45 frames (RF-T-03),
    por lo que en la práctica la forma es \texttt{(45, 75, 4)}. Todos los
    valores deben ser finitos (\texttt{np.isfinite}); cada frame tiene 75
    landmarks con 4 componentes (\texttt{x, y, z, visibility}). \\

    blendshapes & \texttt{float32} & \texttt{(45, 52)} &
    Coeficientes faciales capturados en la misma pasada de MediaPipe que los
    landmarks, con los marcadores no manuales gramaticales de la LSA
    (preguntas, negación, intensidad); la cantidad de filas debe igualar la
    cantidad de frames de landmarks (\texttt{frame\_count}), cada fila debe
    tener exactamente 52 coeficientes (\texttt{BLENDSHAPE\_COUNT}), y todos
    los valores deben ser finitos (validado por
    \texttt{validate\_capture\_blendshapes}). \\
    \bottomrule
    \caption{Diccionario de datos — MuestraDeEntrenamiento (landmarks y blendshapes)}
    \label{tab:dat-dic-muestra}
\end{longtable}

\begin{longtable}{p{2.8cm} p{2.2cm} p{3.2cm} p{5.6cm}}
    \toprule
    \textbf{Atributo} & \textbf{Tipo} & \textbf{Formato / forma} & \textbf{Restricciones} \\
    \midrule
    \endhead
    indice & \texttt{int} & posición en el arreglo & Alineado con la capa de
    salida del modelo cargado (índice 0 = primera clase reconocida). \\

    etiqueta & \texttt{str} & string dentro de un array JSON &
    \texttt{models/custom\_classes.json} es un arreglo JSON de strings,
    escrito por \texttt{train\_lstm.py} antes de iniciar el entrenamiento,
    a partir de los subdirectorios existentes en
    \texttt{dataset/custom/}. \\
    \bottomrule
    \caption{Diccionario de datos — VocabularioDelModelo}
    \label{tab:dat-dic-vocabulario}
\end{longtable}

\begin{longtable}{p{2.8cm} p{2.2cm} p{3.2cm} p{5.6cm}}
    \toprule
    \textbf{Atributo} & \textbf{Tipo} & \textbf{Formato / forma} & \textbf{Restricciones} \\
    \midrule
    \endhead
    archivo & \texttt{str} & nombre de archivo \texttt{.pth} &
    \texttt{models/lstm.pth}: mejor checkpoint del último reentrenamiento. \\

    contenido & binario & \texttt{state\_dict} serializado &
    Escrito con \texttt{torch.save} al finalizar el entrenamiento, sobre el
    checkpoint de mejor exactitud de validación; no es legible sin cargarlo
    con PyTorch. \\

    frames\-Por\-Muestra & \texttt{int} & escalar &
    45 (ventana fija de reconocimiento, ver RF-01.3). \\
    \bottomrule
    \caption{Diccionario de datos — ModeloDeReconocimiento}
    \label{tab:dat-dic-modelo}
\end{longtable}

\begin{longtable}{p{2.8cm} p{2.2cm} p{3.2cm} p{5.6cm}}
    \toprule
    \textbf{Atributo} & \textbf{Tipo} & \textbf{Formato / forma} & \textbf{Restricciones} \\
    \midrule
    \endhead
    modelo & binario & \texttt{.onnx} &
    \texttt{voices/es\_AR-daniela-high.onnx}; activo de terceros
    (voz Piper TTS en español argentino), no generado por LSA Salud. \\

    configuración & JSON & sidecar &
    \texttt{voices/es\_AR-daniela-high.onnx.json}; metadata fonémica y de
    configuración que \texttt{PiperVoice.load} requiere junto al
    \texttt{.onnx}. \\
    \bottomrule
    \caption{Diccionario de datos — ArchivoDeVoz}
    \label{tab:dat-dic-voz}
\end{longtable}

\begin{longtable}{p{2.8cm} p{2.2cm} p{3.2cm} p{5.6cm}}
    \toprule
    \textbf{Atributo} & \textbf{Tipo} & \textbf{Formato / forma} & \textbf{Restricciones} \\
    \midrule
    \endhead
    clave (traducción) & \texttt{str} & texto normalizado &
    Secuencia de palabras normalizada, ordenada y sin duplicados; clave de
    \texttt{\_translation\_cache} (\texttt{OrderedDict}, máximo 200
    entradas, política LRU). \\

    valor (traducción) & \texttt{str} & oración en español &
    Resultado de la traducción (plantilla, caché o modelo de lenguaje). \\

    clave (audio) & \texttt{str} & texto de la oración &
    Clave de \texttt{\_audio\_cache} (\texttt{OrderedDict}, máximo 50
    entradas, política LRU). \\

    valor (audio) & \texttt{str} & WAV codificado en base64 &
    Ninguno de los dos cachés se persiste en disco: viven únicamente en la
    memoria del proceso del backend y se pierden al reiniciarlo. \\
    \bottomrule
    \caption{Diccionario de datos — Cachés en memoria (traducción y audio)}
    \label{tab:dat-dic-caches}
\end{longtable}

\begin{longtable}{p{2.8cm} p{2.2cm} p{3.2cm} p{5.6cm}}
    \toprule
    \textbf{Atributo} & \textbf{Tipo} & \textbf{Formato / forma} & \textbf{Restricciones} \\
    \midrule
    \endhead
    id & \texttt{number} & identificador de mensaje &
    Campo del tipo \texttt{ChatMessage} (\texttt{web/types.ts}). \\

    role & \texttt{str} & \texttt{"PACIENTE"} \textbar\ \texttt{"RECEPCIONISTA"} &
    Indica el emisor del mensaje. \\

    text & \texttt{str} & texto libre & Contenido del mensaje. \\

    time & \texttt{str} & hora de envío & Sin persistencia: el arreglo
    \texttt{chatHistory} vive solo como estado de React de cada pestaña y
    se sincroniza entre pestañas del mismo navegador mediante
    \texttt{Broadcast\allowbreak Channel("nexora-chat")}; no hay backend ni disco
    involucrados. Se pierde al cerrar o recargar la pestaña, y el personal
    de recepción o enfermería lo borra explícitamente en ambas caras de la
    estación al finalizar cada atención (CU-12). \\
    \bottomrule
    \caption{Diccionario de datos — Conversación / Mensaje (no persistido)}
    \label{tab:dat-dic-conversacion}
\end{longtable}

\subsubsection{Entidades de consentimiento, auditoría y reporte (RF-03/04/06)}

Las Tablas~\ref{tab:dat-dic-auditoria} y~\ref{tab:dat-dic-reporte} detallan
estas dos entidades. A diferencia de \textit{MuestraDeEntrenamiento}
(un archivo \texttt{.npy} por muestra), cada una es una \textbf{línea JSON}
dentro de un único archivo compartido y de solo anexado
(\texttt{logs/audit.jsonl}, \texttt{logs/reports.jsonl}): no hay clave
primaria de archivo, sino orden de aparición en el archivo; la escritura
está protegida por un \texttt{threading.Lock} compartido entre ambos
archivos (\texttt{\_append\_jsonl} en \texttt{app.py}).

\begin{longtable}{p{2.8cm} p{2.2cm} p{3.6cm} p{5.2cm}}
    \toprule
    \textbf{Atributo} & \textbf{Tipo} & \textbf{Formato / forma} & \textbf{Restricciones} \\
    \midrule
    \endhead
    tipo & \texttt{str} & \texttt{"prediction"} \textbar\ \texttt{"sentence"} \textbar\ \texttt{"report"} \textbar\ \texttt{"consent"} \textbar\ \texttt{"translation\_displayed"} \textbar\ \texttt{"fin\_atencion"} &
    Identifica el evento; \texttt{"consent"} cubre RF-03.1 (no existe una
    entidad ni un archivo de consentimiento separados) y
    \texttt{"fin\_atencion"} cubre el cierre de atención de CU-12. \\

    timestamp & \texttt{str} & ISO 8601 UTC &
    Generado por \texttt{\_log\_audit} con
    \texttt{datetime.now(timezone.\allowbreak utc).\allowbreak isoformat()};
    no es un nombre de archivo, es un campo más de la línea JSON. \\

    campos según tipo & \texttt{float} \textbar\ \texttt{str} \textbar\ \texttt{int} &
    varía por \texttt{tipo} &
    \texttt{prediction}: \texttt{lstm\_ms}, \texttt{confidence}.
    \texttt{sentence}: \texttt{translation\_ms}, \texttt{tts\_ms},
    \texttt{total\_ms}, \texttt{source}, \texttt{word\_count}.
    \texttt{report}: \texttt{station}.
    \texttt{consent}: \texttt{station}, \texttt{consent}
    (\texttt{"accepted"}\textbar\texttt{"declined"}).
    \texttt{translation\_displayed}: \texttt{station}, \texttt{e2e\_ms}
    (validado finito, $0 \leq$ e2e\_ms $\leq$ 60000, por
    \texttt{validate\_audit\_payload}).
    \texttt{fin\_atencion}: \texttt{station}. \\
    \bottomrule
    \caption{Diccionario de datos — EventoDeAuditoría}
    \label{tab:dat-dic-auditoria}
\end{longtable}

\begin{longtable}{p{2.8cm} p{2.2cm} p{3.6cm} p{5.2cm}}
    \toprule
    \textbf{Atributo} & \textbf{Tipo} & \textbf{Formato / forma} & \textbf{Restricciones} \\
    \midrule
    \endhead
    texto & \texttt{str} & texto libre, máx.\ 500 caracteres &
    Texto del \texttt{ChatMessage} señalado en el momento del reporte
    (validado por \texttt{validate\_report\_payload}); como la conversación
    no se persiste (Tabla~\ref{tab:dat-dic-conversacion}), el reporte
    guarda directamente el texto, no una referencia al \texttt{id} del
    mensaje. \\

    estacion & \texttt{str} & \texttt{"paciente"} \textbar\ \texttt{"recepcion"} &
    Estación desde la que se reportó (derivada de \texttt{usePathname} en
    \texttt{ChatBubble}). \\

    palabras & \texttt{list[str]} \textbar\ \texttt{null} & hasta 20
    palabras, cada una $\leq$ 32 caracteres &
    Palabras crudas reconocidas del mensaje, cuando están disponibles;
    ausente (\texttt{null}) si el mensaje no las tiene. \\

    timestamp & \texttt{str} & ISO 8601 UTC &
    Generado por \texttt{report\_translation} con
    \texttt{datetime.now(timezone.\allowbreak utc).\allowbreak isoformat()}. \\
    \bottomrule
    \caption{Diccionario de datos — ReporteDeTraducción}
    \label{tab:dat-dic-reporte}
\end{longtable}

\subsection{Mapeo físico}
\label{sec:dat-mapeo-fisico}

La Tabla~\ref{tab:dat-mapeo-fisico} resume el mapeo entre cada entidad
lógica y su representación física en disco (o su ausencia de
representación física, para los elementos que LSA Salud mantiene
intencionalmente sin persistir).

\small
\begin{longtable}{>{\raggedright\arraybackslash}p{4.3cm} >{\raggedright\arraybackslash}p{5.6cm} >{\raggedright\arraybackslash}p{3.9cm}}
    \toprule
    \textbf{Entidad lógica} & \textbf{Patrón de ruta física} & \textbf{Formato de archivo} \\
    \midrule
    \endhead
    Muestra\-De\-Entrenamiento (landmarks) &
    \texttt{dataset/\allowbreak custom/\allowbreak<etiqueta>/\allowbreak<timestamp\_ms>.npy} &
    numpy \texttt{float32} \texttt{(45,75,4)} \\

    Muestra\-De\-Entrenamiento (blendshapes) &
    \texttt{dataset/\allowbreak custom/\allowbreak<etiqueta>/\allowbreak<timestamp\_ms>.blendshapes.npy} &
    numpy \texttt{float32} \texttt{(45,52)} \\

    Vocabulario\-Del\-Modelo &
    \texttt{models/custom\_classes.json} &
    JSON --- arreglo de strings \\

    Modelo\-De\-Reconocimiento &
    \texttt{models/lstm.pth} &
    PyTorch \texttt{state\_dict} serializado \\

    Archivo\-De\-Voz &
    \texttt{voices/es\_AR-daniela-high.onnx} +
    \texttt{voices/es\_AR-daniela-high.onnx.json} &
    ONNX + configuración JSON (Piper) \\

    Caché de traducciones &
    en memoria (proceso del backend) &
    no persistido en disco \\

    Caché de audio TTS &
    en memoria (proceso del backend) &
    no persistido en disco \\

    Conversación / Mensaje &
    estado de React en el navegador, sincronizado por
    \texttt{BroadcastChannel} &
    no persistido \\

    Evento\-De\-Auditoría (incluye \texttt{consent}, RF-03.1) &
    \texttt{logs/audit.jsonl} &
    JSON Lines --- un evento anónimo por línea \\

    Reporte\-De\-Traducción &
    \texttt{logs/reports.jsonl} &
    JSON Lines --- un reporte por línea \\
    \bottomrule
    \caption{Mapeo entre entidades lógicas y persistencia física}
    \label{tab:dat-mapeo-fisico}
\end{longtable}

\subsection{Migración futura a base de datos relacional (fuera de alcance)}

El esquema basado en archivos descripto en este entregable cubre las
necesidades de LSA Salud (una instancia local por estación, sin consultas
relacionales, sin historial persistente de conversaciones). Si el proyecto
necesitara en algún momento desplegar múltiples instancias concurrentes
compartiendo el mismo vocabulario y dataset, ofrecer reportes o consultas
sobre las muestras capturadas (por ejemplo, filtrar por rango de fechas o
por desarrollador que capturó cada muestra), o conservar un historial de
conversaciones entre sesiones, esos requisitos sí justificarían introducir
una base de datos relacional: los metadatos de cada muestra
(\textit{etiqueta}, \textit{timestamp}, quién la capturó) y el registro de
\textit{Conversación}/\textit{Mensaje} son los candidatos naturales para
migrar a tablas, mientras que los arreglos \texttt{.npy} y los pesos del
modelo probablemente seguirían viviendo como archivos referenciados desde
esas tablas (un patrón híbrido común para datos binarios de gran tamaño).
Este análisis queda explícitamente fuera del alcance del sistema descripto
en este entregable y no implica ningún compromiso de implementación.
